% TeXplates/Mathematics.tex
% v0.2.1
% https://github.com/HoyanMok/TeXplates
\documentclass[openany]{book} 
% \documentclass{ctexbook} 如果用中文
% \documentclass[10pt,a4paper]{ctexart}  字体大小和纸张大小，默认分别为10pt和letterpaper
% 五号 = 10.5pt，小四=12pt，四号=14pt
% 其他可选参量如twocolumn, 两行排版

\title{Geometry}
\author{Hoyan Mok\thanks{hoyanmok@outlook.com}}
\date{\today} % 自动生成日期

\usepackage{../TeXplatesMathematics}

\addbibresource{Geometry.bib} % 把这里改成实际的文件名


\begin{document}
\frontmatter
\tableofcontents

\mainmatter{}
\part{Basic Geometry}
\chapter{Geometry in Regions of a Space}
\section{Co-odinate and its transformation}
\begin{definition}[Jacobian]
A transformation of co-odinate from $\boldsymbol{x}$ to $\boldsymbol{y}$
\begin{align*}
	\bv y(\bv x ) = y^i(x^j)\basis e_i=y(x^j).
\end{align*}
Its \indexbf{Jacobian}:
\begin{align}\label{Jacobian of transformation}
	\bv J = 
	\begin{pmatrix}
		\cfrac{\partial y^i}{\partial x^j}
	\end{pmatrix}
\end{align}
\end{definition}
A \indexbf{vector} $\bv u$ at point $\bv x_0$ under such transformation would follow:
\begin{align}\label{vector}
	v^i = \left.
			\frac{\partial y^i}{\partial x^j}
		\right|_{\bv x_0} u^i
\end{align}
\begin{align*}
	\text{i.e.}&&
	\bv v = \bv J_0\bv u
\end{align*}

A linear form (\indexbf{covector}) $\ell \colon \bv x \mapsto \ell(\bv x) = l_i x^i $ under such transformation would follow:
\begin{align}\label{linear form}
	l_i' \dif y^i = l_j \dif x^j 
	&&\Rightarrow&&
	l'_i = \left.
			\frac{\partial x^j}{\partial y_i}
		\right|_{\bv x_0} l_j
\end{align}
\begin{align*}
	\text{i.e.}&&
	\bv l' = \bv l \bv J^{-1}_0
\end{align*}

A linear transformation $\mathscr{L} \colon \bv x \mapsto \bv L \bv x$ where $\bv L = ({L^i}_j)_{i, j \in n}$ under such transformation would follow:
\begin{align*}
	\dif y\left(\mathscr L^i(\bv x)\right) 
	&=(L')^i_{\;j} \dif y^j 
	\\
	&=  \left.\frac{\partial y^i}{\partial x^k}\right|_{\bv x_0} \dif \mathscr L^k (\bv x)
	=\left.\frac{\partial y^i}{\partial x^k}\right|_{\bv x_0} L^k_{\;h} \dif x^h
	=\left.\frac{\partial y^i}{\partial x^k}\right|_{\bv x_0} 
	L^k_{\;h} 
	\left.\frac{\partial x^h}{\partial y_j}\right|_{\bv x_0}\dif y^j
\end{align*} 
\begin{align}\label{linear transformation} 
	(L')^i_{\;j} = 
	\left.\frac{\partial y^i}{\partial x^k}\right|_{\bv x_0} 
	L^k_{\;h} 
	\left.\frac{\partial x^h}{\partial y_j}\right|_{\bv x_0} 
	&&\text{or}&&
	\bv L'  = \bv J_0\bv L\bv J^{-1}_0 
\end{align} 

A bilinear form $\mathscr{B} \colon \bv x \mapsto \bv x^{\mathrm{T}} \bv b \bv x = x^{\;i} b_{ij} x^j$: 
\begin{align*}
	b'_{ij}\dif y^{\;i} \dif y^j 
	= b'_{ij} \left.\frac{\partial y^{\,i}}{\partial x^{\;k}}\right|_{\bv x_0}
	\left.\frac{\partial y^j}{\partial x^h}\right|_{\bv x_0}\dif x^{\;k}\dif x^h 
	=b_{kh} \dif x^{\;k}\dif x^h 
	&&\Rightarrow&&
	b'_{ij} =\left. \frac{\partial x^{\;k}}{\partial y^{\;h}}\right|_{\bv x_0}
	b_{kh} 
	\left.\frac{\partial x^h}{\partial y^j}\right|_{\bv x_0}
\end{align*}
\begin{align}\label{bilinear form}
	\text{i.e.} &&
	\bv b' = (\bv J^{-1}_0)^{\mathrm{T}} \bv b \bv J^{-1}_0
\end{align}

\section{Riemannian and Pseudo-Riemannnian Spaces}

\begin{definition}[Riemannian metric]
	A \indexbf{Riemannian metric} $\bv G$ is a smooth, positive-definite quadrutic form defined on a finite-dimensional vector space over $\mathbb R$.
\end{definition}

Given a basis, we usually denote the Riemannian metric by $\indexmath[gij(x)]{g_{ij}(\bv x)}$.

We can define \indexbf{arc length} $\ell$ and \indexbf{inner product} $\langle , \rangle$ in a \indexbf{Riemannian space} (i.e.\ a vector space equiped with a Riemannian metric):
\begin{equation*}
	\ell := \int_{t_1}^{t_2} \sqrt{g_{ij}[\bv x(t)] \diff{x^i}{t} \diff{x^j}{t}} 
	\dif t,
	\qquad
	\langle \bv u, \bv v \rangle := g_{ij} u^i v^j.
\end{equation*}

We can also introduce the following notation: $u_i := g_{ij}u^j$, which means the linear form $\bv v \mapsto u_i v^i$; 
and $\dif \ell^2 = g_{ij} \dif x^i \dif x^j$. 

\begin{definition}[Euclidean metric]
	If a metric $\bv G(\bv x)$ is said to be \emphbf{Euclidean}\index{Euclidean metric} if there exists a coodinates $\bv y(\bv x)$ s.t.\ 
	\begin{equation*}
		g_{ij} = \delta_{k\ell} \pdiff{y^k}{x^i} \pdiff{y^\ell}{x^j}.
	\end{equation*}  

	Such coodinates $\bv y(\bv x)$ are said to be a \indexbf{Euclidean coodinates}.
\end{definition}

\begin{definition}[Pseudo-Riemannian metric]
	A \indexbf{pseudo-Riemannian metric} $\bv G$ is a smooth, indefinite quadrutic form defined on a finite-dimensional vector space over $\mathbb R$.
\end{definition}

A pseudo-Riemannian metric shall have the following cannonical form at some coodinates:
\begin{equation*}
	\bv G = \diag (\eta^2_1, \ldots, \eta^2_p, - \eta^2_{p+1}, \ldots, -\eta^2_n)
\end{equation*}
where $\eta$

By Sylvester's law of inertia, the index of inertia i.e.\ the number of positive terms on the caninical form, shall conserve under any coodinate change.

\begin{definition}[Pseudo-Euclidean metric]
	If a metric $\bv G(\bv x)$ is said to be \emphbf{pseudo-Euclidean}\index{pseudo-Euclidean metric} if there exists a coodinates $\bv y(\bv x)$ s.t.\ 
	\begin{equation*}
		g_{ij} = \sum_{k = 1}^p \pdiff{y^k}{x^i} \pdiff{y^k}{x^j} - \sum_{\ell = p + 1}^n \pdiff{y^\ell}{x^i} \pdiff{y^\ell}{x^j}.
	\end{equation*}  

	Such coodinates $\bv y(\bv x)$ are said to be a \indexbf{pseudo-Euclidean coodinates}.
\end{definition}

We denote a pseudo-Euclidean space by $\mathbb R^n_{p, n-p}$, where $p$ is the index of innertia.
Especially, we call $\mathbb R^4_{1, 3}$ the \indexbf{Minkowski space}, since its significance in relativistic mechanics.

\section{The Simplest Groups of Transformations of a Euclidean Space}

When we say a \indexbf{tansformation} from $\varOmega$ to $\varOmega'$, we refer to a bijective $\varphi$ s.t.\ both $\varphi$ and $\varphi^{-1}$ are smooth, i.e.\ a \indexbf{diffeomorphism}. 
If $\varOmega = \varOmega'$, we would call it the transformation of $\varOmega$.

The transformations form a group, we might call it the \indexbf{transformation group}. 

\begin{definition}[Isometry]
	If a transformation $\varphi$ of $\varOmega$, with a metric $\bv G$, satisfies that $\forall \bv x \in \varOmega$, 
	\begin{equation*}
		g_{ij} = g_{k \ell} \pdiff{\varphi^k}{x^i} \pdiff{\varphi^\ell}{x^j},
	\end{equation*}
	we shall call it a \indexbf{isometry}, or a \indexbf{motion} of the given metric.
\end{definition}

\begin{theorem}[Isometry group]
	The isometries of a region form a group.
\end{theorem}

We might call this group the \indexbf{isometries} or the \indexbf{group of motions}. 
If an isometry preserve the orientation of $\mathbb R^n$, we might call it a \indexbf{proper isometry}. 
The proper isometries is a subgroup of the isometries of $\mathbb R^n$.

\begin{lemma}[Proper isometries of Euclidean space]
	Every proper isometry of $\mathbb R^n$ is either a translation along a vector, or a rotation about a point.
\end{lemma}

The proper isometries can be parameteriesed by a special othogonal matrix $\bv A \in \SO(n)$ (2D special othogonal group), and a vector $\bv v$:
\begin{equation*}
	\bv x \mapsto \bv A \bv x + \bv v.
\end{equation*}

We can find the group a representation by a matrix:
\begin{equation}\label{equation: representation of isometry group}
	\begin{pmatrix}
		\bv A & \bv v \\
		\bv 0^\mathrm T & 1
	\end{pmatrix}
\end{equation}

If $\det \bv A = \pm 1$, then we have the parametrisation of the isometry of the Euclidean space $\mathbb R^n$. 

\begin{theorem}[Mazur-Ulam Theorem]
	The isometry group of $\mathbb R^n$ is the semiproduct of translations ($\simeq \mathbb R^n$) and the othogonal group $\mathrm O(n)$ (rotations and rotatory reflections).
\end{theorem}

The isometry group can be represented by Eq.~\ref{equation: representation of isometry group}.

\begin{theorem}
	$\SO(n)$, as a Lie group, is continuous.
\end{theorem}
\begin{proof}
	find the rotation a basis so that it take canonical form
	\begin{equation*}
		\begin{pmatrix}
			\pm 1 &		   &  &&\\
				& \ddots &  \\
				&		   & \pm 1 \\
			& & & \cos\theta_1 & -\sin\theta_1\\
			& & & \sin\theta_1 & \cos\theta_1\\
			& & &		&				& \ddots & \\
			& & &		&				&		 &  \cos\theta_m & -\sin\theta_m \\
			& & &		&				& 		 &	\sin \theta_m & \cos \theta_m \\
		\end{pmatrix}
	\end{equation*}
	Hence by replacing $\theta_i$ by $t \theta_i$, we can get the rotation from identity continuously by changing $t$ from $0$ to $1$.
\end{proof}

The \indexbf{Galilean group} is the group of \indexbf{Galilean transformation}, characterising by a rotatory reflection $\bv A \in \mathrm O(3)$, a translation $\bv x_0 \in \mathbb R^3$ a velocity $\bv v \in \mathbb R^3$:
\begin{equation*}
	x := (ct, \bv x) \mapsto (ct, \bv A \bv x + \bv x_0 - \bv v t).
\end{equation*}

The Galilean group can be represented as:
\begin{equation}\label{equation: representation of Galilean group}
	\begin{pmatrix}
		\bv A & \bv x_0 & \bv v \\
		\bv 0^\mathrm T & 1 & 0 \\
		\bv 0^\mathrm T & 0 & 1
	\end{pmatrix}	
\end{equation}

\section{Curvature and Torsion}

Let $\gamma \colon [t_1, t_2] \to \mathbb R^3; t \mapsto \bv x(t)$ be a diffeomorphism to a smooth curve in $\mathbb R^3$, the velocity and the acceleration are defined as:
\begin{equation*}
	\bv v := \dot{\bv x} := \diff{\bv x}{t} = \gamma'(t),
	\qquad
	\bv a := \ddot{\bv x} := \diff[2]{\bv x}{t} = \gamma''(t).
\end{equation*}

\begin{definition}[Natural parametre]
	$s(t)$ is the \indexbf{natural parametre} of the curve defined as:
	\begin{equation*}
		s(t) = \int_0^t |\bv v| \dif t.
	\end{equation*}
\end{definition}

We shall denote the velocity and the accelaration by $\bv v_s$ and $\bv a_s$, with natural parametre. 
To change the parametre, we only need:
\begin{equation*}
	\dif s = |\bv v| \dif t,
	\qquad
	\diff{}{s} = \frac{1}{|\bv v|} \diff{}{t}.
\end{equation*}

By differentiating $\bv v_s \cdot \bv v_s = 1$ we know the orthogonality of $\bv a_s$ and $\bv v_s$.

Define three unit vectors for a parameteriesed curve $t \mapsto \bv x(t)$: the tangent vector $\basis t := \bv v/|\bv v|$, the (primary) normal $\basis n := \diff{\basis t}{s} / |\diff{\basis t}{s}|$ and $\basis b = \basis t \times \basis n$. 

$\bv v_s = \basis t$, $\basis n = \basis a_s / |\basis a_s|$.

\begin{definition}[Curvature and torsion]
	The \indexbf{curvature} $\indexmath[kappa]{\kappa}$ is the norm of $\bv a_s$, and the torsion is defined as $\indexmath[tau]{\tau} = \diff{\basis n}{s}$.
\end{definition}

\begin{theorem}[Serret-Frenet formulae for space curve]%
	\index{Serret-Frenet formulae}
	\begin{equation}\label{equation: Serret-Frenet}
		\diff{}{s} \begin{pmatrix}
			\basis t \\ \basis n \\ \basis b
		\end{pmatrix}
		= \begin{pmatrix}
			0 & \kappa & 0 \\
			-\kappa & 0 & \tau \\
			0 & -\tau \\ 0
		\end{pmatrix}
		\begin{pmatrix}
			\basis t \\ \basis n \\ \basis b
		\end{pmatrix}.
	\end{equation}
\end{theorem}

\section{Pseudo-Euclidean Spaces}

$\mathbb R^4_{1, 2}$ (or, more general case, $\mathbb R^n_{1, n-1}$) is called the Minkowski space. 
The metric is always denoted by $\eta_{\mu\nu}$, $\eta_{00} = 1$, $\eta_{ij} = - \delta_{ij}$, $i, j \neq 0$. (By convention, greek index take its value from $0$ to $n-1$)
Considering the physics reality, we usually restrict our study to the region where $x_\mu x^\mu = \eta_{\mu\nu} x^\mu x^\nu \geq 0$, that is, the \indexbf{time-like} region bounded by the \indexbf{light cone} (or \indexbf{isotropic cone}).

A \indexbf{world-line} is the curve $x(t) := (ct, \bv x(t))$, where $t := x^0/c$.
The natural parametre of a world-line is called the \indexbf{proper time} $\indexmath[tau]{\tau}$ (not to be confused with the torsion).

The isomeries of $\mathbb R^4_{1, 3}$ is called the \indexbf{Poincar\'er group}. 

\begin{theorem}
	Any isometry of $\mathbb R^4_{1, 3}$ is an affine transformation.
\end{theorem}
\begin{proof}
	We need to prove that if there points are collinear, the image of them under the isometry is also collinear. 
	After a translation, which is affine, we can set one of the point to be the origin, and the remianing two points are in the same branch (both $t > 0$ or both $t < 0$). 

	\begin{figure}[h]
		\centering
		\includegraphics[width=12cm]{./img/special_reletivity.eps}
		\caption{Pseudo-spheres in $\mathbb R^2_{1,1}$.}
		\label{figure: special relativity}
	\end{figure}

	A pseudo-sphere centred at a point is a hyperboloid (of two sheets) centred at it and the principal axis is along the $x^0$-axis. 
	Since the isometry preserve the distance, under which the points must move along the pseudo-sphere centred at origin.
	As illustrated in Fig.~\ref{figure: special relativity}. 
	If $x_\mu x^\mu = \ell$, and $y^\mu = \lambda x^\mu$, under an isometry, $x \mapsto x'$. 
	The pseudo-sphere centred at $x'$ of radius $(\lambda - 1) \ell$ and thepseudo-sphere centred at origin of radius $\lambda \ell$ are tangent at $y'$, which along with $x'$ and the origin are collinear.	
\end{proof}


\backmatter{}
\nocite{*} % 这个表示列出所有没有在文中被引用的参考文献
\printbibliography[heading=bibliography, title={Bibliography}]
\indexprologue{Here listed the important symbols used in this notes.}
\printindex[symbol]

\printindex
\end{document}